\documentclass[titlepage, fleqn, a4paper, 12pt, twoside]{article}
\usepackage{geometry}
\usepackage{exsheets} %question and solution environments
\usepackage{amsmath, amssymb, amsthm} %standard AMS packages
\usepackage[utf8]{inputenc}
\usepackage{esint} %integral signs
\usepackage{marginnote} %marginnotes
\usepackage{gensymb} %miscellaneous symbols
\usepackage{commath} %differential symbols
\usepackage{xcolor} %colours
\usepackage{cancel} %cancelling terms
\usepackage[free-standing-units,space-before-unit]{siunitx} %formatting units
	\sisetup
	{
		per-mode=fraction,
		fraction-function=\frac
	}
\usepackage{tikz, pgfplots} %diagrams
	\usetikzlibrary{calc, hobby, patterns, intersections, angles, quotes, spy}
	\usetikzlibrary{circuits.logic.US,circuits.logic.IEC}
\usepackage{graphicx} %inserting graphics
\usepackage{hyperref} %hyperlinks
\usepackage{datetime} %date and time
\usepackage{enumerate, enumitem} %numbered lists
\usepackage{float} %inserting floats
\usepackage[american voltages]{circuitikz} %circuit diagrams
\usepackage{setspace} %double spacing
\usepackage{microtype} %micro-typography
\usepackage{listings} %formatting code
	\lstset{language=Matlab}
	\lstdefinestyle{standardMatlab}
	{
		belowcaptionskip=1\baselineskip,
		breaklines=true,
		frame=L,
		xleftmargin=\parindent,
		language=C,
		showstringspaces=false,
		basicstyle=\footnotesize\ttfamily,
		keywordstyle=\bfseries\color{green!40!black},
		commentstyle=\itshape\color{purple!40!black},
		identifierstyle=\color{blue},
		stringstyle=\color{orange},
	}
\usepackage{booktabs}
\usepackage{multirow}
\usepackage{todonotes}
\usepackage[noabbrev,capitalize]{cleveref}
\usepackage[section]{placeins}
\usepackage[style=numeric, backend=biber]{biblatex}
\usepackage{adjustbox}
\usepackage{algpseudocode} %algorithms
\usepackage{algorithm} %algorithms

% \bibliography{<mybibfile>}% ONLY selects .bib file; syntax for version <= 1.1b
\addbibresource{bibliography.bib}% Syntax for version >= 1.2

\newcommand\numberthis{\addtocounter{equation}{1}\tag{\theequation}} %adds numbers to specific equations in non-numbered list of equations

\DeclareMathAlphabet{\mathcal}{OT1}{pzc}{m}{it}

\theoremstyle{definition}
\newtheorem{example}{Example}
\newtheorem{definition}{Definition}

\theoremstyle{theorem}
\newtheorem{theorem}{Theorem}
\newtheorem{law}{Law}

\tikzset{
block/.style = {draw, rectangle, minimum height=3em, minimum width=3em},
tmp/.style  = {coordinate},
sum/.style = {draw, circle, node distance=1cm},
input/.style = {coordinate},
output/.style = {coordinate},
pinstyle/.style = {pin edge={to-,thin,black}}
}

\makeatletter
\@addtoreset{section}{part} %resets section numbers in new part
\makeatother

\newcommand\blfootnote[1]{%
	\begingroup
	\renewcommand\thefootnote{}\footnote{#1}%
	\addtocounter{footnote}{-1}%
	\endgroup
}

\renewcommand{\marginfont}{\scriptsize \color{blue}}

\renewcommand{\tilde}{\widetilde}

\def\doubleunderline#1{\underline{\underline{#1}}}

\SetupExSheets{solution/print = true} %prints all solutions by default

\DeclareMathOperator{\FT}{\mathcal{F}}
\DeclareMathOperator{\IFT}{\mathcal{F}^{-1}}

\DeclareMathOperator{\Q}{\mathcal{Q}}

\newcommand{\NOT}{\overline}

\DeclareMathOperator{\rect}{\mathrm{rect}}
\DeclareMathOperator{\sinc}{\mathrm{sinc}}

%opening
\title{Digital Processing of Single and Multidimensional Signals}
\author{Aakash Jog}
\date{2017-18}

\begin{document}

\pagenumbering{roman}
\begin{titlepage}
\newgeometry{margin=0cm}
\maketitle
\end{titlepage}
\restoregeometry
%\setlength{\mathindent}{0pt}

\blfootnote
{
	\begin{figure}[H]
		\includegraphics[height = 12pt]{cc.pdf}
		\includegraphics[height = 12pt]{by.pdf}
		\includegraphics[height = 12pt]{nc.pdf}
		\includegraphics[height = 12pt]{sa.pdf}
	\end{figure}
	This work is licensed under the Creative Commons Attribution-NonCommercial-ShareAlike 4.0 International License. To view a copy of this license, visit \url{http://creativecommons.org/licenses/by-nc-sa/4.0/}.
} %CC-BY-NC-SA license

\tableofcontents

\clearpage
\section{Lecturer Information}

\textbf{Dr. Raja Giryes}\\
~\\
E-mail: \href{mailto:raja@tauex.tau.ac.il}{raja@tauex.tau.ac.il}\\
~\\

\section{Grading}

\begin{enumerate}
	\item Homework (MATLAB or NumPy): 20\%
	\item Final exam: 20\%
\end{enumerate}

\section{General Structure}

\begin{enumerate}
	\item Multidimensional Fourier analysis
	\item Hilbert spaces and generalized sampling
	\item Spline interpolation
	\item Filter banks and wavelets
	\item Sparsity
	\item Low rank theory or collaborative filtering
	\item Phase retrieval
\end{enumerate}

\clearpage
\pagenumbering{arabic}

\part{Basic Definition and Theorems}

\begin{definition}[Fourier transform]
	For an LTI system, the Fourier transform is defined as
	\begin{align*}
		\FT\left\{ x(t) \right\}j &= \left\langle x,e^{2 \pi j f t} \right\rangle\\
		&= \int\limits_{t \in R^d} x(t) e^{-2 j f^T t} \dif t
	\end{align*}
	\label{def:Fourier_transform}
\end{definition}

\begin{definition}[Tensor product]
	\begin{align*}
		(x_1 \otimes \dots \otimes x_d)(t) &= x_1(t_1) \cdot \dots \cdot x_d(t_d)
	\end{align*}
	for $t \in R^d$.
\end{definition}

\begin{definition}
	The Fourier transform of a tensor product is defined as
	\begin{align*}
		\FT\left\{ x_1 \otimes \dots \otimes x_d \right\}(f) &= \int\limits_{R^d} x_1(t_1) \cdot \dots \cdot x_d(t_d) e^{-2 \pi j (t_1 f_1 + \dots + t_d f_d)} \dif t_1 \dots \dif t_d\\
		&= \left( \FT\left\{ x_1(t) \right\}(f_1) \right) \dots \left( \FT\left\{ x_d(t) \right\}(f_d) \right)
	\end{align*}
\end{definition}

\begin{definition}[$\rect$]
	\begin{align*}
		\rect(t) &=
			\begin{cases}
				0 &;\quad |t| > \frac{1}{2}\\
				1 &;\quad |t| < \frac{1}{2}\\
			\end{cases}
	\end{align*}
\end{definition}

\begin{definition}[$\sinc$]
	\begin{align*}
		\sinc(t) &= \frac{\sin(\pi t)}{\pi t}
	\end{align*}
\end{definition}

\begin{theorem}
	\begin{align*}
		\FT\left\{ \rect(t) \right\} &= \sinc(f)
	\end{align*}
\end{theorem}

\begin{definition}
	A box function is defined to be the product of rect functions in multiple dimensions.
\end{definition}

\begin{theorem}[Shifting in time]
	\begin{align*}
		\FT\left\{ D_p\left( x(t) \right) \right\}(f) &= e^{-2 \pi j f^T p} \FT\left\{ x(t) \right\}(f)
	\end{align*}
	where
	\begin{align*}
		D_p\left( x(t) \right) &= x(t - p)
	\end{align*}
	\label{thm:shifting_in_time}
\end{theorem}

\begin{theorem}[Stretching in time]
	\begin{align*}
		\FT\left\{ S_A\left( x(t) \right) \right\}(f) &= \int\limits_{R^d} x(A t) e^{-j 2 \pi f^T t} \dif t\\
		&= \int\limits_{R^d} x(q) ^{-j 2 \pi f^T A^{-1} q} \frac{\dif q}{|\det A|}\\
		&= \frac{1}{|\det A|} \int\limits_{R^d}  x(q) e^{-j 2 \pi \left( A^{-T} f \right)^T q} \dif q\\
		&= \frac{\FT\left\{ x(t) \right\}\left( A^{-T} f \right)}{|\det A|}\\
		&= \frac{S^{A^{-T}} X(f)}{|\det A|}
	\end{align*}
	where
	\begin{align*}
		S_A\left( x(t) \right) &= x(A t)
	\end{align*}
	where $A$ is a matrix which may represent stretching, rotating, etc.\\
	Hence, if $A$ is equal to a scalar $a$,
	\begin{align*}
		\FT\left\{ x(a t) \right\} &= \frac{\FT\left\{ x(t) \right\}\left( \frac{f}{a} \right)}{|a|}
	\end{align*}
	\label{thm:stretching_in_time}
\end{theorem}

\begin{theorem}[Fourier transform of Gaussian]
	Using \cref{thm:stretching_in_time}, for a single dimensional Gaussian signal,
	\begin{align*}
		\FT\left\{ e^{-t^2} \right\} &= e^{-\pi f^2}
	\end{align*}
	and for a orthogonal multidimensional Gaussian,
	\begin{align*}
		\FT\left\{ e^{-t^T t} \right\} &= e^{-\pi f^T f}
	\end{align*}
	Hence, using \cref{thm:stretching_in_time}, for a general multidimensional Gaussian,
	\begin{align*}
		\FT\left\{ e^{-t C^{-1} t} \right\} &= \frac{1}{\sqrt{\det C}} e^{-\pi f^T C f}
	\end{align*}
	where $C$ is the covariance matrix.
	\label{thm:Fourier_transform_of_Gaussian}
\end{theorem}

\begin{definition}
	A matrix $A$ is said to be unitary if
	\begin{align*}
		A^{-1} &= A^*
	\end{align*}
	where $A^*$ is the conjugate transpose of $A$.
\end{definition}

\begin{theorem}[Rotation]
	Stretching with by a unitary matrix $A$ is equivalent to rotation, and hence denoting $S_A$ by $R_R$,
	\begin{align*}
		\FT\left\{ R_R\left( x(t) \right) \right\} &= R_R \FT\left\{ x(t) \right\}
	\end{align*}
	\label{thm:rotation}
\end{theorem}

\begin{definition}
	A dimension-reducing projection $P$ from $R^d$ to $R^{d - 1}$, is defined as
	\begin{align*}
		P\left( x(t) \right) &= \int\limits_{R^d} x(t_1,\dots,t_d) \dif t_d
	\end{align*}
\end{definition}

\begin{theorem}[Fourier transform of dimension reducing projection]
	The Fourier transform of a dimension reducing projection is
	\begin{align*}
		\FT\left\{ P\left( x(t) \right) \right\} &= \int\limits_{R^{d - 1}} \left( \int\limits_{R^d} x(t_1,\dots,t_d) \dif t_d \right) e^{-j 2 \pi (t_1 f_1 + \dots + t_{d - 1} f_{d - 1})} \dif t_1 \dots \dif t_{d - 1}\\
		&= \int\limits_{R^{d - 1}} \left( \int\limits_{R^d} x(t_1,\dots,t_d) e^{-j 2 \pi t_d (0)} \dif t_d \right) e^{-j 2 \pi (t_1 f_1 + \dots + t_{d - 1} f_{d - 1})} \dif t_1 \dots \dif t_{d - 1}\\
		&= \int\limits_{R^d} x(t) e^{-j 2 \pi f^T t} \dif t \Big|_{f_d = 0}\\
		&= \FT\left\{ x(t) \right\}(f_1,\dots,f_{d - 1},0)\\
		&= \Q\left\{ \FT\left\{ x(t) \right\} \right\}
	\end{align*}
	where $\Q$ is the slicing operator as in \cref{def:slicing_operator}.
	\label{thm:Fourier_transform_of_dimension_reducing_projection}
\end{theorem}

\begin{definition}
	The slicing operator $\Q$ is defined as
	\begin{align*}
		\Q\left\{ X(f_1,\dots,f_d) \right\} &= X\left( f_1,\dots,f_{d - 1},0 \right)
	\end{align*}
	\label{def:slicing_operator}
\end{definition}

\clearpage
\part{Minimum Delay Systems and Group Delay}

\begin{theorem}
	Consider a real system $H(f)$ with input $x(t)$ and output $y(t)$.
	Assume $H(f)$ to be all-pass but with variable phase, i.e.
	\begin{align*}
		\left| H(f) \right| &= 1
	\end{align*}
	If $H(f)$ has a constant delay, i.e.
	\begin{align*}
		H(f) &= e^{-j 2 \pi f \alpha}
	\end{align*}
	then
	\begin{align*}
		y(t) &= x(t - \alpha)
	\end{align*}
\end{theorem}

\begin{definition}[Phase delay]
	The phase delay of a system $H(f)$ for a narrow band input is defined as
	\begin{align*}
		\tau_p &= -\frac{\angle H(f)}{2 \pi f_0}\\
		&= \frac{\varphi_0}{2 \pi f_0}
	\end{align*}
	where $\pm f_0$ are the central frequencies of the narrow band input.
	\label{def:phase_delay}
\end{definition}

\begin{theorem}
	Consider a real system $H(f)$ with input $x(t)$ and output $y(t)$.
	Assume $H(f)$ to be all-pass but with variable phase, i.e.
	\begin{align*}
		\left| H(f) \right| &= 1
	\end{align*}
	Assume
	\begin{align*}
		x(t) &= s(t) \cos\left( 2 \pi f_0 t + \varphi \right)
	\end{align*}
	and equivalently,
	\begin{align*}
		X(f) &= \frac{1}{2} S(f - f_0) e^{j \varphi} + \frac{1}{2} S(f + f_0) e^{-j \varphi}
	\end{align*}
	If $H(f)$ has a phase delay, i.e.
	\begin{align*}
		\angle H(f) &=
			\begin{cases}
				-\varphi_0 &;\quad f > 0\\
				\varphi_0 &;\quad f < 0\\
			\end{cases}
	\end{align*}
	and equivalently,
	\begin{align*}
		H(f) &=
			\begin{cases}
				e^{-j \varphi_0} &;\quad f > 0\\
				e^{j \varphi_0} &;\quad f < 0\\
			\end{cases}
	\end{align*}
	Then,
	\begin{align*}
		Y(f) &= X(f) H(f)\\
		&= \frac{1}{2} S(f - f_0) e^{\varphi - \varphi_0} e^{j (\varphi - \varphi_0)} + \frac{1}{2} S(f + f_0) e^{-j (\varphi - \varphi_0)}
	\end{align*}
	and equivalently,
	\begin{align*}
		y(t) &= s(t) \cos(2 \pi f_0 t + \varphi - \varphi_0)\\
		&= s(t) \cos\left( 2 \pi f_0 (t - \tau_p) + \varphi \right)
	\end{align*}
	where $\tau_p$ is the phase delay as in \cref{def:phase_delay}.
\end{theorem}

\begin{definition}
	The group delay of a system $H(f)$ for a narrow band input is defined as
	\begin{align*}
		\tau_g(f) &= -\dod{}{f} \angle H(f)
	\end{align*}
	\label{def:group_delay}
\end{definition}

\begin{theorem}
	Consider a real system $H(f)$ with input $x(t)$ and output $y(t)$.
	Assume $H(f)$ to be all-pass but with variable phase, i.e.
	\begin{align*}
		\left| H(f) \right| &= 1
	\end{align*}
	Assume
	\begin{align*}
		x(t) &= s(t) \cos\left( 2 \pi f_0 t + \varphi \right)
	\end{align*}
	and equivalently,
	\begin{align*}
		X(f) &= \frac{1}{2} S(f - f_0) e^{j \varphi} + \frac{1}{2} S(f + f_0) e^{-j \varphi}
	\end{align*}
	If $H(f)$ has a non constant phase, such that
	\begin{align*}
		\angle H(\pm f_0) &= \mp \varphi_1
	\end{align*}
	then let $\mp \varphi_0$ be the points of intersection of the tangents to $\angle H(f)$ at $\pm f_0$.\\
	Hence, by Taylor's Formula,
	\begin{align*}
		\angle H(f) &\approx H(f_0) + (f - f_0) \left( \dod{}{f} \angle H(f) \Big|_{f = f_0} \right)\\
		&\approx \angle H(f_0) - (f - f_0) \tau_g(f_0)
	\end{align*}
	where $\tau_g$ is the group delay as in \cref{def:group_delay}.\\
	Similarly,
	\begin{align*}
		\angle H(f) &\approx \angle H(-f_0) - (f + f_0) \tau_g(-f_0)
	\end{align*}
	Hence,
	\begin{align*}
		-\varphi_1 &= \angle H(f_0)\\
		-\varphi_0 &= \angle H(f_0) + f_0 \tau_g(f_0)
	\end{align*}
	Hence, expressing the system as a sum of a system with a constant delay and a system with a constant phase, i.e.
	\begin{align*}
		H(f) &= H_A(f) + H_B(f)
	\end{align*}
	where
	\begin{align*}
		\angle H_A(f) &=
			\begin{cases}
				\varphi_0 &;\quad f < 0\\
				-\varphi_0 &;\quad f > 0\\
			\end{cases}\\
		\angle H_B(f) &= \tau_g(f_0)
	\end{align*}
	the output is
	\begin{align*}
		y(t) &= y_A(t) + y_B(t)\\
		&= s\left( t - \tau_g(f_0) \right) \cos\left( 2 \pi f_0 \left( t - \tau_g(f_0) \right) + \varphi - \varphi_0 \right)\\
		&= s\left( t - \tau_g(f_0) \right) \cos\left( 2 \pi f_0 (t - \tau_g - \tau_p) + \varphi \right)
	\end{align*}
\end{theorem}

\clearpage
\part{Hilbert Spaces}

\section{Inner Product Spaces}

\begin{definition}[Inner product spaces]
	An inner product space is defined to be a space with an inner product such that
	\begin{align*}
		\langle x,y \rangle &= \langle y,x \rangle^*\\
		\langle x , a y + b z \rangle &= a \langle x,y \rangle + b \langle x,z \rangle\\
		\|x\| &= \langle x,x \rangle\\
		&\ge 0
	\end{align*}
	and $\langle x,x \rangle = 0$ iff $x = 0$.
	Also, $\|x - y\|$ is the distance between $x$ and $y$.
\end{definition}

\begin{theorem}
	Let $M$ be an operator from the inner product space $H$ to the inner product space $S$.
	Then,
	\begin{align*}
		\langle M x , y \rangle &= \langle x , M^* y \rangle
	\end{align*}
	for all $x \in H$ and $y \in S$, and where
	\begin{align*}
		M^* &= \overline{M^T}
	\end{align*}
	is the complex transpose of $M$.
\end{theorem}

\begin{definition}[Hilbert space]
	An inner product, normed, complete vector space is said to be a Hilbert space.
\end{definition}

\end{document}
